\chapter{Transformer 实现}
\label{lab:30_transformer}

\noindent\textbf{配套文件：}\path{notebook/30_transformer.ipynb}\par
\noindent\textbf{教材关联：}\bookxrefshort{ch:rev-tokenization}、\bookxrefshort{ch:attention}、\bookxrefshort{ch:position-representation}、\bookxrefshort{ch:rev-transformer}、\bookxrefshort{ch:rev-autoregressive}、\bookxrefshort{ch:long-context}。

\section*{实验目标}
把注意力、掩码和教师强制连接为可解释的序列到序列实验。

\section*{准备及定位}
先阅读本册实验总则，确认Notebook中的依赖与资产单元。原理计算、库接口迁移和真实模型训练可能具有不同资源要求，应分别记录设备、版本及运行范围。

源码定位可从下列函数开始；应结合前置数据与配置单元阅读。
\begin{itemize}
\item \texttt{my\_encode\_content}
\item \texttt{my\_encode}
\item \texttt{my\_batch\_encode}
\end{itemize}

\section*{操作及观察}
\begin{enumerate}
\item 核对ByT5分词接口、源目标边界、decoder输入与标签移位。
\item 逐层记录注意力头、前馈层、残差及交叉注意力的形状，验证掩码可见范围。
\item 保持基线划分与预算，训练翻译闭环并从起始词元自回归生成；区分训练损失和生成评价。
\item 使用Notebook中的参数映射与契约检查对照torch.nn.Transformer；分析迁移MarianMT时架构与制品的变化。
\end{enumerate}

\section*{结果判读及提交}
提交因果性检查、移位样本、损失曲线和生成错误分析。交叉注意力热图只表示内部权重；不能直接称为语言学对齐或因果解释。

提交时附上所执行单元范围、环境与制品身份、原始结果及失败说明。将未运行项目明确列出，不能用Notebook中已有输出替代本次执行证据。

相关方法的原始定义与适用前提见\cite{vaswani2017attention}；本实验的运行结果仍须按实际配置单独验证。
